\NumberedItem{3}
\noindent\hspace*{\parindent}[Es ist]
\[
F_n=(1-ff)^{-\theta}+\frac{\theta(\theta-1)}{n[+1]}(1-ff)^{-\theta-1}ff-\,{.}\,{.}\,{.}.
\]

\indent Schon in der Formel S. 46 [des Handbuchs\textsuperscript{*)}] ist enthalten
\[
F_n=(1-ff)^{-\theta}F\left(\theta,1-\theta,n+1,-\frac{ff}{1-ff}\right)
\]
\noindent und [es] bleibt also nur übrig nachzuweisen, wie viel Bedeutung dieser Formel übrig bleibt, wenn sie divergirt.
\par\smallskip\noindent{\small[S. 290]}\par
\noindent\hspace*{\parindent}Die Ableitung der Gleichung
\[
F(\theta,n+\theta,n+1,ff)=(1-ff)^{-\theta}F\left(\theta,1-\theta,n+1,-\frac{ff}{1-ff}\right)
\]
\noindent geschieht für diesen Zweck am passendsten, indem man die einzelnen Glieder des zweiten Theils nach Potenzen von $ff$ entwickelt und alles, was einerlei Potenz von $ff$ enthält, zusammenfasst, wo dann die Gleichheit mit dem betreffenden Gliede des ersten Theils leicht erkannt wird. Es ist z. B. das Glied des ersten Theils, welches $f^6$ enthält
\[
\frac{\theta{.}\,\theta+1{.}\,\theta+2}{1{.}\,2{.}\,3}{.}\,
\frac{n+\theta{.}\,n+\theta+1{.}\,n+\theta+2}{n+1{.}\,n+2{.}\,n+3}f^6,
\]
\noindent während aus dem zweiten Theil hervorgeht
\par\medskip\noindent\rule{.25\linewidth}{.4pt}\par\smallskip
{\small\noindent *) [Gemeint ist die Formel, oben S. 350, $F(\alpha,\beta,\gamma,x)=(1-x)^{-\alpha}F\left(\alpha,\gamma-\beta,\gamma,-\frac{x}{1-x}\right)\mathpunct{\raisebox{.35ex}{.}}$]\par}
\clearpage\SourceLeaf{364}\ArithmeticHead{357}{R}{ZUR THEORIE DER REIHE $F(\alpha,\beta,\gamma,x)$.\quad IV.}
\[
\begin{aligned}
&\frac{\theta{.}\,\theta+1{.}\,\theta+2}{1{.}\,2{.}\,3}{.}\,f^6\\[.8em]
&-\frac{\theta+1{.}\,\theta+2}{1{.}\,2}{.}\,\frac{\theta{.}\,1-\theta}{1{.}\,n+1}{.}\,f^6\\[.8em]
&+\frac{\theta+2}{1}{.}\,\frac{\theta{.}\,\theta+1}{1{.}\,2}{.}\,\frac{1-\theta{.}\,2-\theta}{n+1{.}\,n+2}{.}\,f^6\\[.8em]
&-\frac{\theta{.}\,\theta+1{.}\,\theta+2}{1{.}\,2{.}\,3}{.}\,\frac{1-\theta{.}\,2-\theta{.}\,3-\theta}{n+1{.}\,n+2{.}\,n+3}{.}\,f^6\\[.8em]
&=\frac{\theta{.}\,\theta+1{.}\,\theta+2}{1{.}\,2{.}\,3}{.}\,f^6{.}\,F(-3,1-\theta,n+1,1)\\[.8em]
&=\frac{\theta{.}\,\theta+1{.}\,\theta+2}{1{.}\,2{.}\,3}f^6{.}\,\frac{\Pi n{.}\,\Pi(n+2+\theta)}{\Pi(n+3)\Pi(n+\theta-1)}\\[.8em]
&=\frac{\theta{.}\,\theta+1{.}\,\theta+2}{1{.}\,2{.}\,3}\,
\frac{n+\theta{.}\,n+\theta+1{.}\,n+\theta+2}{n+1{.}\,n+2{.}\,n+3}[{.}\,f^6]
\end{aligned}
\]
\clearpage
